\section{Tanúsítványkezelés és CA folyamatok}

Ez a fejezet bemutatja a saját CA létrehozását, CSR-ek készítését és tanúsítványok aláírását.

\subsection{CA alapok}
\begin{lstlisting}[language=bash]
# CA private key
openssl genrsa -out ca.key 4096
# Self-signed CA certificate
openssl req -x509 -new -nodes -key ca.key -sha256 -days 3650 -out ca.crt -subj "/C=HU/ST=SomeState/L=City/O=Org/OU=IT/CN=MyCA"
\end{lstlisting}

\subsection{Webszerver CSR és aláírás}
\begin{lstlisting}[language=bash]
# Web private key
openssl genpkey -algorithm RSA -out web.key -pkeyopt rsa_keygen_bits:2048
# Create CSR
openssl req -new -key web.key -out web.csr -subj "/C=HU/ST=SomeState/L=City/O=Org/CN=www.example.com"
# CA signing
openssl x509 -req -in web.csr -CA ca.crt -CAkey ca.key -CAcreateserial -out web.crt -days 825 -sha256
\end{lstlisting}

\subsection{Trust store kezelés}
\begin{itemize}
  \item Linux: \texttt{trust anchor ca.crt} vagy rendszerszintű \texttt{/usr/local/share/ca-certificates/} és \texttt{update-ca-certificates}
  \item Browser/OS: tenyérnyi különbségek, dokumentáld a célplatformot
\end{itemize}
